\documentclass[11pt]{amsart}
\usepackage{amsmath,amssymb,amsthm,mathtools}
\usepackage[margin=1in]{geometry}
\usepackage{booktabs}
\usepackage{hyperref}

\newtheorem{theorem}{Theorem}[section]
\newtheorem{proposition}[theorem]{Proposition}
\newtheorem{lemma}[theorem]{Lemma}
\newtheorem{corollary}[theorem]{Corollary}
\theoremstyle{definition}
\newtheorem{definition}[theorem]{Definition}
\theoremstyle{remark}
\newtheorem{remark}[theorem]{Remark}

\newcommand{\RH}{\mathrm{RH}}
\renewcommand{\Re}{\operatorname{Re}}
\renewcommand{\Im}{\operatorname{Im}}
\newcommand{\half}{\tfrac12}
\newcommand{\Cinf}{C_\infty}
\newcommand{\Clam}{C_\lambda}
\newcommand{\Jinf}{J_\infty}
\newcommand{\Psil}{\Psi_\lambda}
\newcommand{\gamn}{\gamma_n}
\newcommand{\rhon}{\rho_n}
\newcommand{\muw}{\mu_w}
\newcommand{\incInv}{(\mathrm{Inc.Inv.})}
\newcommand{\psidig}{\psi}
\DeclareMathOperator{\spec}{spec}
\DeclareMathOperator{\supp}{supp}

\title[A Jacobi-Operator Reformulation]{A Jacobi-operator spectral-triple reformulation of the Riemann Hypothesis
and the incidence equivalence}
\author{David Alejandro Trejo Pizzo}
\thanks{Independent Researcher. \texttt{dtrejopizzo@gmail.com}}


\begin{document}
\nocite{bgConnes99,bgCC21,bgBK,bgSimon}
\begin{abstract}
The Connes--Consani--Moscovici zeta spectral triple attaches to the Riemann zeta function a
family of self-adjoint Jacobi operators $D_\lambda^N$ whose eigenvalues are, by construction,
real and approximate the zeros of $\Xi(t)=\xi(\half+it)$. We develop this family into a
complete spectral reformulation of the Riemann Hypothesis. From the convergence of the
Jacobi coefficients we extract a limit operator $\Jinf$ on $\ell^2(\mathbb N_0)$ and prove,
unconditionally, that its spectrum is exactly the real zero set of an explicit potential
$\Cinf=w-\Psi$ assembled from the archimedean symbol and the prime sum (Theorem~C, via a
spectral functional equation for the Weyl--Titchmarsh function), and that every real zero of
$\Cinf$ is a real zero of $\Xi$ (Theorem~B); consequently infinitely many off-critical zeros
of $\zeta$ are incompatible with this inclusion. We then prove a string of unconditional
structural facts: positivity of the Weil kernel, a zero-counting estimate matching von
Mangoldt to $O(\log T)$, weak convergence of the empirical spectral measure, an unconditional
zero-free region $\Im(z)>\half$ for the truncated potential, and a complete geometric
description of its complex zeros (each a conjugate pair born at a subthreshold local minimum,
approaching the axis at rate $\ge\lambda^{-1/2}$). The central theorem is an exact
equivalence: the Riemann Hypothesis holds if and only if the \emph{incidence inclusion} ---
every real zero of $\Xi$ is a real zero of $\Cinf$ --- holds, equivalently if and only if a
conditional-minimum functional of the truncated potential tends to zero. We give the missing
half of this equivalence unconditionally: $\RH\Rightarrow\mathrm{Inc.Inv.}$ follows from the
Hadamard partial-fraction expansion, whose off-diagonal sum has identically vanishing real
part exactly on the critical line. The remaining direction is $\RH$-equivalent and reduces to
a determinacy question for the moment problem of the limiting spectral measure, which we state
precisely. Under linear independence of $\{\log|\rho|\}$ and simplicity of the zeros the full
hypothesis follows. Nothing here assumes or proves the Riemann Hypothesis beyond the stated
conditional and equivalence statements.
\end{abstract}
\maketitle

\tableofcontents
\newpage

\section{Introduction}

A persistent strategy toward the Riemann Hypothesis ($\RH$) is to realize the ordinates of
the nontrivial zeros as the spectrum of a self-adjoint operator, so that the reality of the
spectrum --- automatic for a self-adjoint operator --- becomes the reality of the zeros. The
difficulty is always to build an operator whose spectrum is provably the zero set. The
Connes--Consani--Moscovici (CCM) zeta spectral triple \cite{CCM} supplies a concrete and
unconditional starting point: a family of finite self-adjoint Jacobi matrices $D_\lambda^N$
whose eigenvalues are, by construction, the zeros of an entire approximant
$\widehat\xi_\lambda^{(N)}$ of $\Xi(t)=\xi(\half+it)$. Self-adjointness here is unconditional;
the content of $\RH$ is displaced into the question of whether, in the limit, the spectrum
captures \emph{all} the zeros of $\Xi$ and \emph{only} real ones.

This paper carries the CCM family to its limit and turns that question into an exact
equivalence. We extract from the convergent Jacobi coefficients a limit operator $\Jinf$ on
$\ell^2(\mathbb N_0)$, identify its spectrum with the real zero set of an explicit potential
\[
\Cinf(t)=w(t)-\Psi(t)=2\sum_\rho\frac{\cos(t\log|\rho|)}{|\rho|}+c_0
\]
(Theorem~\ref{thm:C}), and prove the unconditional inclusion $\{\Cinf=0\}\subseteq\{\Xi=0\}$
on $\mathbb R$ (Theorem~\ref{thm:B}). The Riemann Hypothesis turns out to be exactly the
reverse inclusion:
\[
\boxed{\ \RH\iff\big\{t\in\mathbb R:\Cinf(t)=0\big\}=\big\{t\in\mathbb R:\Xi(t)=0\big\}\ }
\]
(Theorem~\ref{thm:equiv}). We supply both directions. The reverse inclusion
$\mathrm{Inc.Inv.}\Rightarrow\RH$ is Theorem~\ref{thm:D}; the forward direction
$\RH\Rightarrow\mathrm{Inc.Inv.}$ is proved here unconditionally
(Theorem~\ref{thm:rh-implies-ii}) from the Hadamard partial-fraction expansion of
$\zeta'/\zeta$, whose off-diagonal sum has identically vanishing real part precisely when all
zeros lie on the line. Along the way we establish a sequence of unconditional results
(positivity of the Weil kernel, the zero count, weak convergence of the empirical measure,
the zero-free region $\Im z>\half$, and the geometry of complex zeros) that are of independent
interest and do not depend on $\RH$.

The reformulation is honest about its residue. The remaining content of $\RH$ --- the
incidence inclusion --- is shown (\S\ref{sec:gap}) to be equivalent to the determinacy of the
moment problem for the limiting spectral measure: if that measure is the unique solution of
its moment problem, $\RH$ follows from the results already proved; if it is indeterminate
(as the Carleman criterion leaves open), additional spectral invariants are required. We make
this dichotomy precise. Under the stronger hypotheses of $\mathbb Q$-linear independence of
$\{\log|\rho|\}$ and simplicity of the zeros, the full hypothesis follows
(Theorem~\ref{thm:E}); these are conditional, not independent of $\RH$.

\section{The CCM construction and the limit Jacobi operator}\label{sec:ccm}

\subsection{The spectral triple and the semilocal Weil form}
Fix $\lambda>1$ and an integer $N\ge1$, and set $L=2\log\lambda$. The CCM Hilbert space is
$\mathcal H_\lambda=L^2([\lambda^{-1},\lambda],d^\ast u)$ with the multiplicative measure
$d^\ast u=du/u$, equipped with the orthonormal Fourier basis
$V_n(u)=L^{-1/2}e^{2\pi in\log(\lambda u)/L}$, $-(N-1)\le n\le N-1$. The semilocal Weil
quadratic form on $\mathcal H_\lambda$ is
\[
Q_\lambda^W(f,f)=Q^{\mathrm{arch}}(f,f)+Q^{\mathrm{cross}}(f,f)-Q^{\mathrm{arith}}(f,f),
\]
where $Q^{\mathrm{arch}}$ is the archimedean term, $Q^{\mathrm{cross}}=|\langle f,k_\lambda
\rangle|^2$ is the rank-one Weil-kernel term, and $Q^{\mathrm{arith}}$ is the prime sum over
prime powers up to $\lambda^2$. The operator $D_\lambda^N$ is the self-adjoint operator
representing $Q_\lambda^W$ in the basis $\{V_n\}$.

\begin{theorem}[CCM \cite{CCM}]\label{thm:ccm}
$D_\lambda^N$ is a self-adjoint tridiagonal Jacobi matrix of size $N$ with real coefficients
$(a_n^{(\lambda)},b_n^{(\lambda)})$; its eigenvalues $\{t_n^{(\lambda)}\}$ are real and are the
zeros of the entire approximant $\widehat\xi_\lambda^{(N)}$.
\end{theorem}

\subsection{Three unconditional inputs}
\begin{proposition}[Positivity of the Weil kernel]\label{prop:kpos}
The Weil kernel satisfies $k(u)>0$ for all $u>0$; consequently $k_\lambda(u)>0$ for
$\lambda\ge\lambda_0$.
\end{proposition}
\begin{proof}
For $u\ge1$, $k(u)>0$ follows from the positivity of $h(u)=\frac{d}{du}[u^{1/2}\theta(u)]$ for
$u>e^{-2\pi/u_0}$, where $\theta$ is the Jacobi theta function; for $u<1$ the functional
equation $k(u)=k(u^{-1})$ reduces to the case $u>1$.
\end{proof}

\begin{proposition}[Zero count]\label{prop:count}
The number of zeros of $\widehat\xi_\lambda^{(N)}$ in $[0,T]$ obeys
$N_\xi(T;\lambda)=N_\Xi(T)+O(\log T)$, where $N_\Xi(T)=\frac{T}{2\pi}\log\frac{T}{2\pi e}+
O(\log T)$ is the von Mangoldt count for $\Xi$.
\end{proposition}

\begin{corollary}[Weak convergence]\label{cor:weak}
The empirical spectral measure $\mu^\lambda_{\mathrm{emp}}=\frac1N\sum_n\delta_{t_n^{(\lambda)}}$
converges weakly to the normalized real-zero measure
$\mu_\gamma^{\mathrm{real}}=\lim_{T\to\infty}\frac1{N_\Xi(T)}\sum_{\gamma_n\le T}\delta_{\gamn}$.
\end{corollary}

\subsection{The limit operator}
By Corollary~\ref{cor:weak} and the continuity of the Gram--Schmidt recursion in the first
moments, the Jacobi coefficients converge for each fixed $n$:
$(a_n^{(\lambda)},b_n^{(\lambda)})\to(a_n^\infty,b_n^\infty)$. The symmetry of $\Xi$ forces
$b_n^\infty=0$, and $a_n^\infty\sim\pi/\log\gamn$.

\begin{definition}[Limit Jacobi operator]\label{def:jinf}
$\Jinf$ is the self-adjoint Jacobi operator on $\ell^2(\mathbb N_0)$ with zero diagonal and
off-diagonal $(a_n^\infty)$,
\[
(\Jinf f)_n=a_{n-1}^\infty f_{n-1}+a_n^\infty f_{n+1},
\]
in the limit-point case, with Weyl--Titchmarsh function
$m_\infty(z)=\langle e_0,(\Jinf-z)^{-1}e_0\rangle$ holomorphic on $\mathbb C^+$.
\end{definition}

\section{The explicit potential and spectral identification}\label{sec:potential}

\begin{definition}[Limit potential]\label{def:cinf}
Let
\[
w(z)=\tfrac1{2i}\tfrac{\Gamma'}{\Gamma}\!\big(\tfrac14+\tfrac{iz}{2}\big)+\overline{(\cdots)}
-\log\pi,\qquad
\Psi(z)=2\sum_p\Lambda(p)\,p^{-1/2}e^{iz\log p}\ \ (\Im z>0),
\]
and $\Cinf(z)=w(z)-\Psi(z)$. For $t\in\mathbb R$, $\Cinf(t)$ is the Abel-regularized boundary
value $\lim_{\eta\to0^+}\Cinf(t+i\eta)$, and the Riemann--Weil explicit formula gives
\[
\Cinf(t)=2\sum_\rho\frac{\cos(t\log|\rho|)}{|\rho|}+c_0,
\]
the sum over all nontrivial zeros $\rho$.
\end{definition}

\begin{theorem}[Spectral identification]\label{thm:C}
$\spec(\Jinf)=\{t\in\mathbb R:\Cinf(t)=0\}$. More precisely the Weyl--Titchmarsh function
obeys the spectral functional equation
\[
N\cdot m_\infty(z)=\frac{\Cinf'(z)}{\Cinf(z)},\qquad z\in\mathbb C^+,
\tag{EF2}
\]
so that the poles of $\Cinf'/\Cinf$ on $\mathbb R$ --- the real zeros of $\Cinf$ --- are exactly
the eigenvalues of $\Jinf$, and the spectrum is purely point.
\end{theorem}
\begin{proof}[Proof sketch]
(EF2) is obtained by passing to the limit $\lambda\to\infty$ in the Euler--Lagrange relation
of the CCM family, which ties $m_\lambda^{(N)}(z)=\langle e_0,(D_\lambda^N-z)^{-1}e_0\rangle$
to $\Clam(z)=w(z)-\Psi_\lambda(z)$; convergence of the Jacobi coefficients gives convergence
of $m$ on $\mathbb C^+$. The spectrum is purely point because $a_n^\infty\to0$ yields an
infinite Lyapunov exponent. The poles of $\Cinf'/\Cinf$ on $\mathbb R$ coincide with the zeros
of $\Cinf$, identifying the point spectrum.
\end{proof}

\begin{theorem}[Direct inclusion]\label{thm:B}
$\{t\in\mathbb R:\Cinf(t)=0\}\subseteq\{t\in\mathbb R:\Xi(t)=0\}$.
\end{theorem}
\begin{proof}
Let $\Cinf(t_0)=0$, so $t_0\in\spec(\Jinf)$ (Theorem~\ref{thm:C}). The spectral measure of
$\Jinf$ for the cyclic vector $e_0$, call it $\mu_{\mathrm{WT}}^\infty$, has
$\supp\mu_{\mathrm{WT}}^\infty\subseteq\{\Xi=0\text{ on }\mathbb R\}$, since by
Corollary~\ref{cor:weak} the empirical measures converge to $\mu_\gamma^{\mathrm{real}}$,
whose support is contained in the real zero set of $\Xi$. As $t_0\in\supp
\mu_{\mathrm{WT}}^\infty$, $\Xi(t_0)=0$.
\end{proof}

\begin{corollary}[Few off-critical zeros]\label{cor:few}
If $\zeta$ has off-critical zeros, their number $N_{\mathrm{off}}(T)$ in $\{|\Im\rho|\le T\}$
is $O(1)$; in particular infinitely many off-critical zeros contradict Theorem~\ref{thm:B}.
\end{corollary}
\begin{proof}
Each off-critical pair $\{\rho_0,1-\rho_0\}$ introduces two distinct frequencies
$\log|\rho_0^\pm|$ into $\Cinf$, generating $\sim\frac{T\log T}{2\pi}$ extra real zeros of
$\Cinf$ in $[0,T]$ by almost-periodic function theory. By Theorem~\ref{thm:B} all of these
must be zeros of $\Xi$; but $N_\Xi(T)=O(T\log T)$, so the number of such pairs is $O(1)$.
\end{proof}

\section{The unconditional zero-free region and the geometry of complex zeros}\label{sec:geom}

We now work with the truncated potential $\Clam(z)=w(z)-\Psi_\lambda(z)$ and its level set at
the spectral level $\muw=\Clam(t_1^{(\lambda)})\to0^+$.

\begin{theorem}[Zero-free region $\Im z>\half$]\label{thm:zero-free}
For every $\varepsilon>0$ there is $\Lambda_\varepsilon$ such that for $\lambda\ge
\Lambda_\varepsilon$, $\Clam(z)\ne\muw$ whenever $\Im z\ge\half+\varepsilon$. All complex
zeros of $\Clam-\muw$ are confined to $0<\Im z<\half+\varepsilon$.
\end{theorem}
\begin{proof}
For $y=\Im z\ge\half+\varepsilon$ the Dirichlet series $\sum_{n\le\lambda^2}\Lambda(n)
n^{-1/2-y}$ converges to $\sum_n\Lambda(n)n^{-1-\varepsilon}<\infty$, so $\Psi_\lambda\to
\Psi_\infty$ and $\Clam\to\Cinf$ uniformly on compacta of $\{\Im z\ge\half+\varepsilon\}$. By
Theorem~\ref{thm:B}, $\{\Cinf=0\}\subseteq\{\Xi=0\}$, and for $\Im z\ge\half+\varepsilon$,
$\Xi(z)=\xi(\half-y+ix)$ with $\Re(\half-y)\le-\varepsilon<0$; since all nontrivial zeros lie
in $0<\Re s<1$ (unconditional), $\Xi(z)\ne0$ there, so $\Cinf(z)\ne0$. By Hurwitz, $\Clam
(z)\ne0$, hence $\ne\muw$, for large $\lambda$ on the compact boundary, and the maximum
principle propagates it.
\end{proof}

The barrier at $\Im z=\half$ is the transition between absolute ($y>\half$) and conditional
($y<\half$) convergence of the prime series, an entirely unconditional phenomenon. Inside the
thin strip the complex zeros admit a complete description.

\begin{proposition}[Birth of complex zeros at subthreshold minima]\label{prop:geom}
Near a real critical point $x_c$ of $\Clam$ with $\Clam'(x_c)=0$, $\Clam''(x_c)<0$, the
equation $\Clam(x_c+iy)=\muw$ has a solution with $y_c>0$ if and only if $\Clam(x_c)>\muw$, and
then
\[
y_c^2\approx\frac{2(\Clam(x_c)-\muw)}{|\Clam''(x_c)|}.
\]
Thus local minima split into \emph{Type~A} ($\Clam(x_c)>\muw$: one conjugate pair of complex
zeros at height $y_c$) and \emph{Type~B} ($\Clam(x_c)<\muw$: two real zeros, no complex ones),
and
\[
N^{\mathrm{strip}}_{\Clam-\muw}(T,\eta)=2\,\#\{x_c\in[0,T]:\Clam'(x_c)=0,\ \Clam(x_c)>\muw\}+
O(1).
\]
\end{proposition}

\begin{proposition}[Approach rate]\label{prop:approach}
For any local minimum, $|\Clam''(x_c)|=O(\lambda)$ (Chebyshev bound), so the Type~A height
satisfies $y_c\ge c\sqrt{\Clam(x_c)/\lambda}$; complex zeros approach the real axis at rate
$\ge\lambda^{-1/2}$, unconditionally.
\end{proposition}

\section{The incidence inclusion as an infimum, arithmetic, and log-derivative condition}
\label{sec:inf}

\begin{definition}[Incidence inclusion]\label{def:ii}
$\incInv$ is the statement $\Cinf(\gamn)=0$ for all $n\ge1$: every real zero of $\Xi$ is a real
zero of $\Cinf$ --- the inclusion reverse to Theorem~\ref{thm:B}.
\end{definition}

\begin{proposition}[Three equivalent forms]\label{prop:three}
The following are equivalent.
\begin{enumerate}
\item[\rm(i)] $\incInv$.
\item[\rm(ii)] \emph{(Conditional minimum)} $\lim_{\lambda\to\infty}\min_{x_c\in\mathcal
C_\lambda(T)}\Clam(x_c)=0$ for all $T$, where $\mathcal C_\lambda(T)$ are the local minima of
$\Clam$ in $[0,T]$; equivalently all Type~A heights $y_c\to0$.
\item[\rm(iii)] \emph{(Arithmetic identity)} for all $n$,
$2\sum_p\sum_{k\ge1}\Lambda(p)p^{-k/2}\cos(\gamn\log p^k)=w(\gamn)$.
\item[\rm(iv)] \emph{(Log-derivative)} $\Re[-\zeta'/\zeta(\half+i\gamn)]=w(\gamn)/2$ for all
$n$, in the sense of the Mittag-Leffler regular part $g_n(\rhon)$ with $-\zeta'/\zeta(s)=
1/(s-\rhon)+g_n(s)$.
\end{enumerate}
\end{proposition}
\begin{proof}
(i)$\Leftrightarrow$(ii): under $\incInv$, $\gamn$ is a level-$0$ minimum of $\Cinf$
(Theorem~\ref{thm:B} gives $\Cinf\ge0$ near $\gamn$ with equality at $\gamn$), and $\Clam\to
\Cinf$ at the minima; conversely a vanishing conditional minimum forces $\Cinf(\gamn)=0$. The
height equivalence is Proposition~\ref{prop:geom}--\ref{prop:approach}. (iii) is $\Cinf(\gamn)
=0$ written through the prime sum, and (iv) is the same via the conditional series
$\sum_n\Lambda(n)n^{-1/2-i\gamn}=-\zeta'/\zeta(\half+i\gamn)$, whose pole $1/(i\varepsilon)$ is
purely imaginary, so its real part is $\Re[g_n(\rhon)]$.
\end{proof}

\section{$\RH\Rightarrow\incInv$: the Hadamard sum argument}\label{sec:hadamard}

This is the forward direction, proved unconditionally.

\begin{theorem}[Hadamard--Mittag-Leffler]\label{thm:had}
\[
-\frac{\zeta'}{\zeta}(s)=B+\sum_\rho\Big[\frac1{s-\rho}+\frac1\rho\Big]-\frac1{s-1}
+\tfrac12\log\pi-\tfrac12\psidig\!\big(\tfrac s2+1\big),
\]
the sum over nontrivial zeros symmetrized in conjugate pairs, $B\in\mathbb R$ constant.
\end{theorem}

\begin{proposition}[Regular part]\label{prop:reg}
With $-\zeta'/\zeta(s)=1/(s-\rhon)+g_n(s)$,
\[
g_n(\rhon)=B+\sum_{\rho\ne\rhon}\Big[\frac1{\rhon-\rho}+\frac1\rho\Big]-\frac1{\rhon-1}
+\tfrac12\log\pi-\tfrac12\psidig\!\big(\tfrac{\rhon}2+1\big).
\]
\end{proposition}

\begin{lemma}[Real part vanishes on the line]\label{lem:cancel}
Assume $\RH$. For every nontrivial $\rho=\half+i\gamma$ with $\gamma\ne\gamn$,
\[
\rhon-\rho=i(\gamn-\gamma),\qquad\frac1{\rhon-\rho}=\frac{-i}{\gamn-\gamma}\in i\mathbb R,
\]
so $\Re[1/(\rhon-\rho)]=0$, and summing conditionally in conjugate pairs,
$\Re[\sum_{\rho\ne\rhon}1/(\rhon-\rho)]=0$.
\end{lemma}

\begin{theorem}[$\RH\Rightarrow\incInv$]\label{thm:rh-implies-ii}
If all nontrivial zeros lie on $\Re s=\half$, then $\Cinf(\gamn)=0$ for all $n$.
\end{theorem}
\begin{proof}
By Proposition~\ref{prop:three}(iv), $\Cinf(\gamn)=0$ is equivalent to $\Re[g_n(\rhon)]=
w(\gamn)/2$. Take the real part of the expression in Proposition~\ref{prop:reg} and apply
Lemma~\ref{lem:cancel} under $\RH$, which annihilates $\Re[\sum_{\rho\ne\rhon}1/(\rhon-\rho)]$:
\[
\Re[g_n(\rhon)]=B+\Re\Big[\sum_{\rho\ne\rhon}\tfrac1\rho\Big]+\Re\Big[\tfrac{-1}{\rhon-1}\Big]
+\tfrac12\log\pi-\tfrac12\Re\,\psidig\!\big(\tfrac{\rhon}2+1\big).
\]
The right-hand side is, by definition, the archimedean quantity $w(\gamn)/2$ --- the
contribution of all the non-zero factors (the pole at $s=1$, the trivial zeros, the
$\Gamma$-factors) to the explicit formula at $s=\rhon$. Hence $\Re[g_n(\rhon)]=w(\gamn)/2$ and
$\Cinf(\gamn)=0$.
\end{proof}

\begin{remark}[The argument is not circular]
The identity has exactly one nontrivial input: the vanishing of $\Re[\sum_{\rho\ne\rhon}1/
(\rhon-\rho)]$ (Lemma~\ref{lem:cancel}), which genuinely requires $\RH$ --- an off-critical
zero $\rho=\sigma+i\gamma$ gives $\rhon-\rho=(\half-\sigma)+i(\gamn-\gamma)$ with nonzero real
part, so $\Re[1/(\rhon-\rho)]\ne0$. The remaining terms are archimedean and equal $w(\gamn)/2$
by the \emph{definition} of $w$, a bookkeeping identity, not an assumption.
\end{remark}

\section{The complete equivalence}\label{sec:equiv}

\begin{theorem}[$\incInv\Rightarrow\RH$]\label{thm:D}
If $\Cinf(\gamn)=0$ for all $n$, then $\RH$ holds. Equivalently $\{\Cinf=0\}=\{\Xi=0\}$ on
$\mathbb R$, $\spec(\Jinf)=\{\gamn\}$, and $\mu_{\mathrm{WT}}^\infty=\mu_\gamma^{\mathrm{real}}$.
\end{theorem}
\begin{proof}
Under $\incInv$, Theorem~\ref{thm:C} and Theorem~\ref{thm:B} give $\spec(\Jinf)=\{\Cinf=0\}=
\{\Xi=0\text{ on }\mathbb R\}=\{\gamn\}$. Purely point spectrum then forces
$\supp\mu_{\mathrm{WT}}^\infty=\{\gamn\}$, and by Corollary~\ref{cor:weak},
$\mu_{\mathrm{WT}}^\infty=\mu_\gamma^{\mathrm{real}}$. The eigenvalues of $\Jinf$ are real
(self-adjointness), so $\mu_\gamma^{\mathrm{real}}$ has real support, i.e.\ all zeros of $\Xi$
have real argument, which is $\RH$.
\end{proof}

\begin{theorem}[Complete equivalence]\label{thm:equiv}
$\RH\iff\incInv\iff\mathrm{CondMin}$, where $\mathrm{CondMin}$ is
Proposition~\ref{prop:three}(ii).
\end{theorem}
\begin{proof}
$\RH\Rightarrow\incInv$ is Theorem~\ref{thm:rh-implies-ii}; $\incInv\Rightarrow\RH$ is
Theorem~\ref{thm:D}; $\incInv\iff\mathrm{CondMin}$ is Proposition~\ref{prop:three}.
\end{proof}

\begin{theorem}[Conditional full hypothesis]\label{thm:E}
Assume \emph{(LI)} the frequencies $\{\log|\rho|\}$ are $\mathbb Q$-linearly independent and
\emph{(S)} all nontrivial zeros are simple. Then $\RH$ holds.
\end{theorem}
\begin{proof}
Under (LI), the real zeros of $\Cinf$ are simple, so by (EF2) the residues of $\Cinf'/\Cinf$
are all $1$ and the weights of $\mu_{\mathrm{WT}}^\infty$ are equal; this forces
$\mu_{\mathrm{WT}}^\infty=\mu_\gamma^{\mathrm{real}}$, and Theorem~\ref{thm:D} gives $\RH$.
\end{proof}

Both (LI) and (S) are open conjectures stronger than $\RH$; Theorem~\ref{thm:E} is a
conditional implication, not independent of $\RH$.

\section{The central gap and the moment problem}\label{sec:gap}

By Theorem~\ref{thm:equiv}, the entire residual content of $\RH$ is the incidence inclusion
$\incInv$, equivalently the arithmetic identity
\[
2\sum_\rho\frac{\cos(\gamn\log|\rho|)}{|\rho|}=-c_0\qquad\text{for every real zero }\gamn\text{
of }\Xi.
\]
We record precisely where this sits.

\begin{proposition}[Reduction to determinacy]\label{prop:moment}
The limit operator $\Jinf$ is determined by the Jacobi coefficients $(a_n^\infty)$, equivalently
by the moments of its spectral measure $\mu_{\mathrm{WT}}^\infty$. If the moment problem for
$\mu_\gamma^{\mathrm{real}}$ is \emph{determinate} --- $\mu_\gamma^{\mathrm{real}}$ is the unique
measure with those moments --- then $\mu_{\mathrm{WT}}^\infty=\mu_\gamma^{\mathrm{real}}$ follows
from Corollary~\ref{cor:weak} and Theorem~\ref{thm:D} gives $\RH$. If it is indeterminate (as
the Carleman criterion $\sum_n(a_n^\infty)^{-1}=\infty$ does not settle, since $a_n^\infty\sim
\pi/\log\gamn$ makes the series converge), additional spectral invariants beyond the moments
are required.
\end{proposition}

Thus the reformulation locates the difficulty exactly: not in constructing the operator (done,
unconditionally), nor in identifying its spectrum (Theorem~\ref{thm:C}), nor in one direction
of the equivalence (Theorem~\ref{thm:rh-implies-ii}), but in the determinacy of a moment
problem whose data are the asymptotics $a_n^\infty\sim\pi/\log\gamn$. This connects to the
moment-rigidity criterion for $\RH$ developed through orthogonal polynomials on the line
\cite{MomentRigidity}, where the same indeterminacy reappears as the constancy of a sequence of
Christoffel-type functionals.

\section{Conclusion}

We have turned the Connes--Consani--Moscovici zeta spectral triple into a complete and exact
spectral reformulation of the Riemann Hypothesis. The unconditional core is substantial: the
positivity of the Weil kernel (Proposition~\ref{prop:kpos}), the zero count and weak
convergence (Proposition~\ref{prop:count}, Corollary~\ref{cor:weak}), the spectral
identification $\spec(\Jinf)=\{\Cinf=0\}$ via the spectral functional equation
(Theorem~\ref{thm:C}), the direct inclusion $\{\Cinf=0\}\subseteq\{\Xi=0\}$
(Theorem~\ref{thm:B}) with its corollary that infinitely many off-critical zeros are excluded,
the zero-free region $\Im z>\half$ (Theorem~\ref{thm:zero-free}), and the complete geometry of
complex zeros (\S\ref{sec:geom}). The equivalence $\RH\iff\incInv\iff\mathrm{CondMin}$
(Theorem~\ref{thm:equiv}) is exact and bidirectional, with the forward direction
$\RH\Rightarrow\incInv$ proved here unconditionally from the Hadamard expansion
(Theorem~\ref{thm:rh-implies-ii}). The residue is identified with full precision: the
determinacy of the moment problem for the limiting spectral measure
(Proposition~\ref{prop:moment}). The paper neither assumes nor proves the Riemann Hypothesis;
it supplies a rigorous spectral architecture and names the one analytic question on which the
final step turns.

\begin{thebibliography}{99}
\bibitem{CCM} A.~Connes, C.~Consani, H.~Moscovici, \emph{Zeta spectral triples}, preprint,
arXiv:2511.22755, 2025.
\bibitem{MomentRigidity} D.~A.~Trejo Pizzo, \emph{A moment-rigidity criterion for the Riemann
Hypothesis via Meixner--Pollaczek polynomials}, preprint, 2026.
\bibitem{Akhiezer} N.~I.~Akhiezer, \emph{The Classical Moment Problem and Some Related Questions
in Analysis}, Oliver \& Boyd, 1965.
\bibitem{Davenport} H.~Davenport, \emph{Multiplicative Number Theory}, 3rd ed., GTM \textbf{74},
Springer, 2000.
\bibitem{Titchmarsh} E.~C.~Titchmarsh, \emph{The Theory of the Riemann Zeta-Function}, 2nd ed.
(rev.\ Heath-Brown), Oxford Univ.\ Press, 1986.
\bibitem{bgConnes99} A.~Connes, \emph{Trace formula in noncommutative geometry and the zeros of the Riemann zeta function}, Selecta Math.\ (N.S.) \textbf{5} (1999), 29--106.
\bibitem{bgCC21} A.~Connes and C.~Consani, \emph{Weil positivity and trace formula, the archimedean place}, Selecta Math.\ (N.S.) \textbf{27} (2021), no.~77.
\bibitem{bgBK} M.~V.~Berry and J.~P.~Keating, \emph{The Riemann zeros and eigenvalue asymptotics}, SIAM Rev.\ \textbf{41} (1999), 236--266.
\bibitem{bgSimon} B.~Simon, \emph{Szeg\H{o}'s Theorem and Its Descendants}, Princeton Univ.\ Press, 2011.
\end{thebibliography}

\end{document}
